% ******************************************************** %
% DOCUMENT INFORMATION                                     %
% ******************************************************** %
%                                                          %
% Purpose of Document                                      %
% -------------------                                      %
% Some Specific Remarks on LaTeX                           %
%                                                          %
% Institution                                              %
% -----------                                              %
% University of Applied Sciences and Arts Northwestern     %
% Switzerland, School of Engineering                       %
%                                                          %
% Degree Program                                           %
% --------------                                           %
% Electrical Engineering and Information Technology, BSc.  %
%                                                          %
% Author & Copyright                                       %
% ------------------                                       %
% Raphael Frey, raphael.frey@students.fhnw.ch              %
%               rmfrey@runbox.com                          %
%                                                          %
% Date: 2017-APR-26                                        %
% ******************************************************** %

\documentclass{beamer}                % Presentation Version
%\documentclass[trans]{beamer}         % Transparency Version
%\documentclass[handout]{beamer}       % Handout Version

% ------------------------------------------------------------ Beamer Setup %<<<
%\addtobeamertemplate{background canvas}{\transdissolve[duration=1]\hspace{-0.29em}}{}
%\usetheme[titleformat=smallcaps,numbering=none]{metropolis}
\usetheme[titleformat=smallcaps,progressbar=frametitle]{metropolis}
%>>>

% ---------------------------------------------------------------- Packages %<<<
\usepackage[ngerman]{babel}
\usepackage{xcolor-solarized}
\usetikzlibrary{calc}
\usetikzlibrary{positioning}
\usetikzlibrary{arrows}
\usetikzlibrary{calc}
\usetikzlibrary{fit}
\usetikzlibrary{spy}
\usetikzlibrary{backgrounds}
\usetikzlibrary{shapes.symbols}
\usepackage{dirtree}
\usepackage{minted}
\setminted{%
    %style=xcode,
    %style=trac,
    %style=paraiso-light,
    %style=lovelace,
    style=murphy,
    bgcolor=solarized-base3,
    linenos=true,
}
\usepackage{booktabs}
%\setbeamercolor{alerted text}{fg=solarized-red}
%>>>


% ------------------------------------------------------------------ Macros %<<<
% Arrow for use in text
\def\tikzrarrow{%
    \tikz[baseline=-0.67ex]\draw[double,very thick,-{stealth}](0,0)--(0.5,0);}
\newlength{\yOffs}
\newlength{\xOffs}
\newcommand*\code[1]{\texttt{#1}}
\newcommand\source[1]{%
    \tikz[remember picture,overlay] 
        \node[align=left,anchor=south west,text=solarized-base1,font=\tiny]
        at (current page.south west)
        {Quelle: #1};
}
\newcommand\contrastnode[1]{%
    \tikz[
        baseline=(n0.base),
        rounded corners=2mm,
        inner sep=5pt,
        text=solarized-base2,
    ] 
        \node[fill=solarized-base01] (n0) {#1}; 
}
%>>>

% ------------------------------------------------------------- Title Setup %<<<
\title{\vspace*{4em}\Huge\LaTeX}
\subtitle{\hfill Trouble on the Horizon}
%\date{\today}
\date{26. April 2017}
\author{%
    Raphael Frey%
    \hfill%
    \scriptsize\texttt{%
        \href{mailto:rmfrey@runbox.com}%
        {rmfrey@runbox.com}}}

%\institute{IME}
\titlegraphic{\includegraphics[height=2em]{images/fhnwLogoDE-solarized-base02.eps}}
\subject{LaTeX -- Kurze Einf\"uhrung}
\keywords{LaTeX Introduction Primer FHNW Overview}
%http://web.fhnw.ch/cd/corporate-design/logos-fur-die-hochschulen
%>>>

% **************************************************************************** %
\begin{document}                                                              
% **************************************************************************** %

\frame[plain]{\titlepage} % -------------------------------------- TITLE FRAME %

\begin{frame} % ------------------------------------- NICE-TO-HAVE PACKAGES %<<<
    \frametitle{Lohnenswerte Packages}
    \begin{itemize}
        \item hyperref
        \item microtype
        \item babel/polyglossia
        \item inputenc
        \item amsmath
        \item fontenc, fontspec, luatextra: https://tex.stackexchange.com/a/44701/131649
    \end{itemize}
\end{frame}%>>>


\begin{frame} % ---------------------------------------------- BIBTEX INTRO %<<<
    \frametitle{Bibliographie -- Was, Wozu?}
    \centering
    \setlength\tabcolsep{1pt}
    \renewcommand\arraystretch{1.5}
    \begin{tabular}{p{20mm}rclp{34mm}}
        \toprule
        \scshape Layer    & \multicolumn{3}{l}{\scshape Komponente} & \scshape Zweck \\
        \midrule
        \LaTeX            & \contrastnode{biblatex}          & \contrastnode{natbib} & \contrastnode{native} & Bereitstellen von \mintinline{tex}|\cite|  and friends \\
        externer Arbeiter & \contrastnode{\scshape Bib\TeX}  & \contrastnode{biber}  &        & Verkn\"upfen von \code{.tex}- und \code{.bib}-Dateien  \\
        Datenbank-Datei   & \contrastnode{\code{.bib}}       & \contrastnode{andere} &        & Referenzdaten speichern                                \\
        \bottomrule
    \end{tabular}
    \source{https://tex.stackexchange.com/a/299286/131649}
\end{frame}%>>>

\begin{frame} % ------------------------------------------BIBTEX FLOW CHART %<<<
    \frametitle{Bib\TeX\ -- Was, Woher, Wohin?}
    \hspace*{-1em}%
    \begin{tikzpicture}[
        draw=solarized-base02,
        text=solarized-base02,
        rounded corners=2mm,
        every node/.append style={
            draw,
            thick,
            fill=solarized-base2,
            % technically not necessary with the global rounded corners option:
            rounded corners=2mm, 
            inner sep=2mm,
            font=\fontsize{9}{10}\selectfont,
        },
        inputfile/.style={
            draw=solarized-orange,
            double,
        },
        signal/.append style={
            text=solarized-base2,
            fill=solarized-magenta,
            rounded corners=1mm,
        },
    ]
        \setlength{\xOffs}{5mm}
        \setlength{\yOffs}{6mm}

        \node[inputfile] (tex) 
            {\code{.tex}};

        \node (latex1)
            [xshift=\xOffs,yshift=-\yOffs,rotate=-45,signal,shape = signal,signal from =west,signal to=east,below right of=tex]
            {latex};

        \node (aux1) 
            [xshift=2*\xOffs,yshift=-\yOffs,right of=latex1]
            {\code{.aux}};

        \node (latex2)
            [xshift=2.5*\xOffs,signal,shape = signal,signal from =west,signal to=east,right of=aux1]
            {latex};

        \node (bibtex)
            [yshift=-2.33 * \yOffs,signal,shape = signal,signal from =west,signal to=east,below of=aux1]
            {bibtex};

        \node (aux2) 
            [xshift=2.5*\xOffs,right of=latex2]
            {\code{.aux}};

        \node (latex3)
            [xshift=2*\xOffs,rotate=-45,signal,shape = signal,signal from =west,signal to=east,right of=aux2]
            {latex};

        \node (bbl)
            [yshift=-\yOffs,below of=latex2]
            {\code{.bbl}};

        \node[inputfile] (bib)
            [xshift=-9* \xOffs,yshift=-\yOffs,below of=latex2]
            {\code{.bib}};

        \node[inputfile] (bst)
            [yshift=-\yOffs,below of=bib]
            {\code{.bst}};

        \node (blg)
            [yshift=-\yOffs,below of=bbl]
            {\code{.blg}};

        \node (pdf) 
            [xshift=8.5 * \xOffs,right of=blg]
            {\color{solarized-red}\code{.pdf,.dvi}};

        \draw[thick,-{stealth}]
            (tex.south) to[out=-90,in=135] (latex1.west);

        \draw[thick,-{stealth}]
            (tex.east) to[out=0,in=180] (latex2.west);

        \draw[thick,-{stealth}]
            (tex.east) to[out=0,in=135] (latex3.west);

        \draw[thick,-{stealth}]
            (latex1.east) to[out=-45,in=180] (aux1.west);

        \draw[thick,-{stealth}]
            (aux1.east) -- (latex2.west);

        \draw[thick,-{stealth}]
            (latex2.east) -- (aux2.west);

        \draw[thick,-{stealth}]
            (latex3.east) to[out=-45,in=90] (pdf.north);

        \draw[thick,-{stealth}]
            (bib.east) to[out=0,in=180] (bibtex.west);

        \draw[thick,-{stealth}]
            (bst.east) to[out=0,in=180] (bibtex.west);

        \draw[thick,-{stealth}]
            (bibtex.east) to[out=0,in=180] (bbl.west);

        \draw[thick,-{stealth}]
            (bbl.north) -- (latex2.south);

        \draw[thick,-{stealth}]
            (bibtex.east) to[out=0,in=180] (blg.west);

        \draw[thick,-{stealth}]
            (aux2.east) to[out=0,in=135] (latex3.west);

        \draw[thick,-{stealth}]
            (aux1.south) -- (bibtex);

    \end{tikzpicture}

    \source{Mittelbach et al.,\emph{The \LaTeX\ Companion}, 2nd Ed., p.688, Fig. 12.1}
\end{frame}%>>>

\begin{frame}[fragile] % ------------------------ NATIVE BIBLIOGRAPHY: CODE %<<<
    \frametitle{Bibliography -- Native}

    \inputminted{tex}{code/native.tex}
\end{frame}
%>>>

\begin{frame} % ------------------------------- NATIVE BIBLIOGRAPHY: RESULT %<<<
    \frametitle{Bibliography -- Native}

    \fbox{\includegraphics[width=\textwidth,clip,viewport=45mm 220mm 170mm 257mm]{code/native.pdf}}
\end{frame}%>>>

\begin{frame}[fragile] % ------------------------------------- BIBTEX: CODE %<<<
    \frametitle{Bib\TeX}

    \code{.tex}-File:

    \inputminted{tex}{code/bibtex.tex}
\end{frame}
%>>>

\begin{frame}[fragile] % --------------------------- BIBTEX: REFERENCES.BIB %<<<
    \frametitle{Bib\TeX}

    \code{references.bib}-File:

    \inputminted{tex}{code/references.bib}
\end{frame}
%>>>

\begin{frame} % -------------------------------------------- BIBTEX: RESULT %<<<
    \frametitle{Bib\TeX}

    \fbox{\includegraphics[width=\textwidth,clip,viewport=45mm 220mm 170mm 257mm]{code/bibtex.pdf}}
\end{frame}%>>>

\begin{frame} % ------------------------------------------- BIBTEX: CAVEATS %<<<
    \frametitle{Bib\TeX\ IEEEtran}

    \begin{itemize}
        \item
            \alert{Vorsicht:} Es    gibt     auch    eine    \textsf{IEEEtran}
            \LaTeX-Klasse! (\code{IEEEtran.cls})
        \item
            Wir   sind  aber   nur  am   \textsc{Bib}\TeX-Style  interessiert:
            \code{IEEEtran.\alert{bst}}
        \item
            Dokumentation:         \code{IEEEtran\_HOWTO.pdf}        respektive
            \code{IEEEtran\_bst\_HOWTO.pdf}
        \item
            \href{http://ctan.org/tex-archive/macros/latex/contrib/IEEEtran/bibtex}
                 {\nolinkurl{ctan.org/tex-archive/macros/latex/contrib/IEEEtran/bibtex}}
        \item
            \href{http://www.ieee.org/documents/IEEEtranBST.zip}
                 {\nolinkurl{ieee.org/documents/IEEEtranBST.zip}}
        \item
            \href{https://github.com/alpenwasser/TeX/tree/master/bibliography/official-IEEE-resources}
                 {\nolinkurl{https://github.com/alpenwasser/TeX/bibliography/official-IEEE-resources}}
    \end{itemize}
\end{frame}%>>>

\begin{frame} % ------------------------------------------- BIBTEX: CAVEATS %<<<
    \frametitle{Bib\TeX\ -- Die h\"assliche Seite}

    \begin{itemize}
        \item
            Referenzieren   und  Erstellen   der  Bibliographie   funktioniert
            $\pm$ automatisch, aber
        \item
            es    gibt    16    verschiedene   IEEE    citation    types    in
            \code{IEEEtran.bst}, und
        \item
            welcher Typ  unter welchen  Umst\"anden wie  auszuf\"ullen
            ist, ist eine (stratosph\"arisch hohe) Kunst f\"ur sich.
        \item
            Meinungsverschiedenheiten sind  vorprogrammiert, selbst  wenn alle
            beteiligten  Parteien   nach  bestem  Wissen  und   Gewissen  nach
            offiziellen Style Guides vorgehen.
    \end{itemize}

    \tikzrarrow\ Allenfalls mal im Voraus die Bibliographie absegnen lassen.
\end{frame}%>>>

\begin{frame} % -------------------- BIBTEX: ACCESSE DATE, NUMBERED SECTION %<<<
    \frametitle{Bib\TeX\ -- Zugriffsdatum f\"ur Onlinequellen}
    \begin{itemize}
        \item
            Wunsch  FHNW: Bibliographie  als   nummerierte  Section  nach  den
            \emph{Schlussfolgerungen},   \emph{Fazit}    o.\"a.: Aufgabe   des
            \code{.cls}-Files, momentan aber noch nicht implementiert.
        \item
            Zugriffsdatum  f\"ur  URLs: Wird   von  \code{IEEEtran.bst}  nicht
            unterst\"utzt  (implizite  Annahme:  alle URLs  sind  g\"ultig  am
            Publikationsdatum).
            \begin{itemize}
                \item
                    Entweder globale Anmkerung diesbez\"uglich machen, oder
                \item
                    manuell das \code{.bbl}-File editieren.
            \end{itemize}
    \end{itemize}
\end{frame}%>>>

\begin{frame} % ----------------------------------------- MODULAR DOCUMENTS %<<<
    \frametitle{Modulare Dokumente}
    master/root document!
\end{frame}%>>>

\begin{frame} % -----------------------------------------------FLOATS INTRO %<<<
    \frametitle{Floats -- \"Ubersicht}
\end{frame}%>>>

\begin{frame} % -------------------------------------------------- HYPERREF %<<<
    \frametitle{Hyperref}
\end{frame}%>>>

%\begin{frame} % -------------------------------------------------------- FRAME %
%    \frametitle{Einfache Macros}
%\end{frame}

%\begin{frame} % -------------------------------------------------------- FRAME %
%    \frametitle{Modulare Dokumente}
%\end{frame}
%
%\begin{frame} % -------------------------------------------------------- FRAME %
%    \frametitle{Farben}
%\end{frame}
%
%\begin{frame} % -------------------------------------------------------- FRAME %
%    \frametitle{\texttt{minipages}}
%\end{frame}
%
%\begin{frame} % -------------------------------------------------------- FRAME %
%    \frametitle{Code Listings}
%\end{frame}
%
%\begin{frame} % -------------------------------------------------------- FRAME %
%    \frametitle{Externe PDF-Dokumente}
%\end{frame}
%
%\begin{frame} % -------------------------------------------------------- FRAME %
%    \frametitle{Elektrische Schaltungen zeichnen}
%\end{frame}
%
%\begin{frame} % -------------------------------------------------------- FRAME %
%    \frametitle{Daten Plotten}
%    \begin{itemize}
%        \item matlab2tikz
%        \item pgfplots
%    \end{itemize}
%\end{frame}
%
%\begin{frame} % -------------------------------------------------------- FRAME %
%    \frametitle{Alternative Klassen}
%\end{frame}
%
%\begin{frame} % -------------------------------------------------------- FRAME %
%    \frametitle{Support}
%
%    \begin{itemize}
%        \item
%            e-Mail:
%            \texttt{\href{mailto:raphael.frey@students.fhnw.ch}
%            {raphael.frey@students.fhnw.ch}}
%        \item
%            Raum 4.223
%        \item
%            \alert{\href{https://github.com/alpenwasser/TeX/}
%                        {https://github.com/alpenwasser/TeX/}}
%        \item
%            Weiteres Programm: 26. April
%            \begin{itemize}
%                \item
%                    modulare Dokumentstruktur
%                \item
%                    Bibliographie
%                \item
%                    Floats (Tabellen, Bilder, \ldots )
%                \item
%                    PGF/Ti\emph{k}Z
%                %\item
%                %    minipages
%                %\item
%                %    ISO-31
%                \item
%                    Listings, minted, \ldots
%                \item
%                    etc. etc. etc.
%            \end{itemize}
%    \end{itemize}
%\end{frame}

% ---------------------------------------------------------------------------- %
\section<handout:0| trans:0>*{Fragen?} % ------------------------------- FRAME %
% ---------------------------------------------------------------------------- %
\end{document}
% vim: foldenable foldcolumn=4 foldmethod=marker foldmarker=<<<,>>>
